\documentclass[11pt,a4paper]{article}
\usepackage{curso}
\begin{document}
\portada{Lección}{7.8}{Complejidad comparada}

\noindent\textbf{Unidad 7.}\quad\textbf{Objetivos:} comparar pasos cuánticos, de difusión y digitales según el recurso físico que consumen; estimar recursos para
$n=2048$ bits; contrastar HHL con el álgebra lineal termodinámica.
\textbf{Prerrequisitos:} 6.1, 6.4, 7.1--7.7.

\section{Tres clases de paso}
\begin{center}\begin{tabular}{lll}\toprule
paso & número & costo físico\\\midrule
compuerta cuántica & $O(n^3)$ & $\sim2n$ qubits lógicos coherentes, corrección de errores\\
paso de difusión & $O(n^2)$ & $\ge r$ nodos, lectura con $\sim2n$ bits analógicos\\
operación digital & subexponencial (NFS) & tiempo\\\bottomrule\end{tabular}\end{center}

\section{Un número de 2048 bits \evE}
Shor: $\approx4099$ qubits lógicos, $n^3\approx8{,}6\times10^9$ compuertas (aritmética de escuela). Difusión: $\ell_0=4(M+1)(n+2)\approx1{,}68\times10^7$ pasos, sobre hasta
$2^{2048}$ nodos con $\approx4096$ bits de precisión. Criba de cuerpos de números: $\exp\big((64/9)^{1/3}(\ln N)^{1/3}(\ln\ln N)^{2/3}\big)\approx2^{117}$ (sin el
término $o(1)$). Rho de Pollard: $N^{1/4}=2^{512}$. La dificultad se traslada: coherencia (Shor), área y precisión (difusión), tiempo (digital).

\section{HHL frente a álgebra lineal termodinámica}
HHL: $\tilde O(\log(d)\,s^2\kappa^2/\varepsilon)$ para matrices $s$-dispersas con acceso por oráculo; entrega el estado $|x\rangle$, y leer todas las entradas
anula la ventaja. TLA (Unidad 6): entrega el vector clásico, $O(d\kappa^2\varepsilon^{-2})$ en unidades de $\tau_r$, y paga la carga $O(d^2)$. Distintas promesas,
distintas letras pequeñas.

\section{Cuatro preguntas}
(1) ¿Qué es un paso y cuánto cuesta físicamente? (2) ¿Cómo crecen área, energía y precisión? (3) ¿Qué se entrega y cuánto cuesta leerlo? (4) ¿Contra qué línea base
clásica y en qué condiciones?

\begin{aplicacion}
Evaluar afirmaciones con modelos explícitos y etiquetas \evM\ \evS\ \evP\ \evE\ (lección 5.8). El proyecto final incluye una tabla de recursos.
\end{aplicacion}

\begin{error}
\begin{itemize}[nosep]
\item Comparar números de pasos de modelos distintos. \quad Ignorar los $\Omega(r)$ nodos de la difusión.
\item Ignorar el oráculo y la salida cuántica de HHL. \quad Confundir qubits físicos y lógicos.
\end{itemize}
\end{error}

\begin{thebibliography}{9}\small
\bibitem{hhl} A. W. Harrow, A. Hassidim y S. Lloyd, Quantum algorithm for linear systems of equations, \emph{Phys. Rev. Lett.} 103, 150502 (2009).
\bibitem{aaronson} S. Aaronson, Read the fine print, \emph{Nature Physics} 11, 291 (2015).
\bibitem{tratado} Cadavid et al., Tratado de factorización por difusión (data/ del curso), secciones 1.3 y 20, 2026.
\end{thebibliography}
\end{document}
